%%%%%%%%%%%%%%%%%%%%%%%%
%
% $Autor: Gautam Ramesh$
% $Datum: 2025-10-14$
% $Kapitel: ListOfParts$
%
%%%%%%%%%%%%%%%%%%%%%%%%

\chapter{List Of Parts}

Selecting and documenting the correct hardware components is crucial to ensuring reliable and repeatable results in embedded AI systems.  
This section provides the full bill of materials (BOM) used in the assembly of the license plate detection prototype.  
The listed components were chosen based on criteria such as power efficiency, availability, community support, and proven use in prior edge-AI research \cite{nvidia-jetson, edge-ai-survey, embedded-vision, yolov8-performance}.

% -------------------------------------------------------
% Table 1: Detailed Component List
% -------------------------------------------------------
\begin{table}[H]
	\centering
	\caption{Detailed Component List and Functions}
	\begin{tabular}{|L{3.2cm}|L{6.4cm}|L{4.0cm}|L{1.2cm}|}
		\hline
		\textbf{Part Name} & \textbf{Technical Specification} & \textbf{Function in System} & \textbf{Qty} \\ \hline
		Jetson Nano (4\,GB, B01) & Quad-core ARM Cortex-A57, 128-core Maxwell GPU; 40-pin GPIO & Runs YOLOv8 detection + EasyOCR; central compute & 1 \\ \hline
		microSD Card & \SI{32}{GB} UHS-I, Class 10; JetPack~5.1.2 image & OS boot / model and data storage & 1 \\ \hline
		USB Camera & UVC compliant; 1080p@30\,FPS (e.g., LifeCam Show) & Captures live plate images & 1 \\ \hline
		HDMI Display & 1080p monitor with HDMI input & Live visualization and debugging & 1 \\ \hline
		USB Keyboard + Mouse & Wired or 2.4\,GHz dongle & Local setup and control & 1 set \\ \hline
		DC Power Supply & \SI{5}{V}, \SI{4}{A} regulated, barrel jack & Stable power for sustained inference & 1 \\ \hline
		HDMI Cable & High-speed, \(\approx\)1\,m & Connects Nano to display & 1 \\ \hline
		USB Cables & USB-A to device (camera/peripherals) & Camera/peripheral connectivity & 1--2 \\ \hline
		Enclosure (optional) & Jetson-compatible case or reComputer shell & Physical protection and airflow & 1 \\ \hline
		GPIO Jumper Set (optional) & Male--Female, \(\approx\)20\,cm & Future sensors / I/O prototyping & as req. \\ \hline
	\end{tabular}
\end{table}

The above table summarizes the core and supporting components required for the embedded vision pipeline.  
Each item was validated under load to ensure sustained inference performance at approximately 10–15~FPS.  
Using the official \SI{5}{V}/\SI{4}{A} NVIDIA adapter prevents under-voltage throttling of the GPU, as highlighted in \cite{edge-power-stability}.  
Cable lengths are deliberately kept below 1.2\,m to reduce data latency and power loss.

\vspace{1em}

\begin{table}[H]
	\centering
	\caption{Optional Accessories and Spares}
	\begin{tabular}{|L{4.0cm}|L{7.2cm}|L{3.0cm}|}
		\hline
		\textbf{Accessory} & \textbf{When Useful} & \textbf{Notes} \\ \hline
		Powered USB Hub & When multiple peripherals exceed onboard current limit & Reduces system stress on 5V rail \\ \hline
		USB Fan / Heatsink & Continuous inference runs $>$30\,min & Improves thermal margin up to 15\% \cite{thermal-design} \\ \hline
		Tripod / Camera Mount & Required for fixed surveillance perspective & Maintains plate alignment consistency \\ \hline
		Spare microSD & Quick OS re-image / recovery point & Ensures reproducibility of environment \\ \hline
		Cable Ties / Labels & Recommended for lab safety & Prevents mechanical strain and confusion \\ \hline
	\end{tabular}
\end{table}

Accessories in Table~2 are recommended for robustness and testing consistency.  
The fan or heatsink improves long-term thermal stability, which directly impacts inference accuracy on edge devices \cite{thermal-design}.  
Maintaining fixed camera geometry via a tripod minimizes variation in dataset capture.

% -------------------------------------------------------
% TikZ Figure: System Assembly Layout
% -------------------------------------------------------
\begin{figure}[H]
	\centering
	\begin{tikzpicture}[
		every node/.style={font=\scriptsize, align=center},
		core/.style={draw, fill=blue!10, thick, rounded corners=3pt, minimum width=3.2cm, minimum height=1.1cm},
		input/.style={draw, fill=green!10, thick, rounded corners=3pt, minimum width=2.9cm, minimum height=1.1cm},
		output/.style={draw, fill=purple!10, thick, rounded corners=3pt, minimum width=2.9cm, minimum height=1.1cm},
		power/.style={draw, fill=red!10, thick, rounded corners=3pt, minimum width=2.8cm, minimum height=1.1cm},
		storage/.style={draw, fill=orange!10, thick, rounded corners=3pt, minimum width=2.8cm, minimum height=1.1cm},
		optional/.style={draw, fill=gray!15, thick, rounded corners=3pt, minimum width=2.9cm, minimum height=1.1cm},
		arr/.style={-{Latex[length=2mm]}, thick}
		]
		\node[core] (nano) at (0,0) {\textbf{Jetson Nano}\\(AI Compute Core)};
		\node[input, left=2.8cm of nano, yshift=1.2cm] (camera) {USB Camera\\(Image Source)};
		\node[input, left=2.8cm of nano, yshift=-1.2cm] (kbd) {Keyboard / Mouse\\(User Control)};
		\node[output, right=2.8cm of nano, yshift=1.2cm] (monitor) {HDMI Monitor\\(Visualization)};
		\node[storage, right=2.8cm of nano, yshift=-1.2cm] (sd) {microSD Card\\(Storage)};
		\node[power, below=2.4cm of nano] (psu) {5V / 4A Power Adapter};
		\node[optional, above=2.2cm of nano] (gpio) {GPIO Header\\(Sensors / Relays)};
		\draw[arr] (camera.east) -- node[above]{USB 3.0} (nano.west);
		\draw[arr] (kbd.east) -- node[above]{USB 2.0} (nano.west);
		\draw[arr] (nano.east) -- node[above]{HDMI} (monitor.west);
		\draw[arr] (nano.east) -- ++(0.3,-1.1) -- (sd.west) node[midway, below]{I/O Bus};
		\draw[arr] (psu.north) -- node[right]{DC Input} (nano.south);
		\draw[arr, dashed] (gpio.south) -- node[right]{GPIO Pins (3.3V Logic)} (nano.north);
	\end{tikzpicture}
	\caption{System assembly layout showing core Jetson Nano connections with color-coded layers for input, output, power, and optional expansion.}
\end{figure}

This diagram demonstrates the modular architecture of the embedded system.  
Each component is color-coded by function, enabling easy identification during assembly and inspection.  
Such layered representations are commonly used in embedded AI system documentation \cite{system-diagrams, edge-ai-survey}.  
The dashed GPIO link emphasizes scalability, supporting future hardware upgrades such as relay control, vehicle sensors, or IoT integration.

% -------------------------------------------------------
% TikZ Figure: Cable Pack and Ports (quick visual)
% -------------------------------------------------------
\begin{figure}[H]
	\centering
	\begin{tikzpicture}[
		every node/.style={font=\scriptsize, align=center},
		box/.style={draw, rounded corners=2pt, fill=gray!10, minimum width=1.8cm, minimum height=0.8cm},
		port/.style={draw, rounded corners=2pt, fill=blue!10, minimum width=1.8cm, minimum height=0.8cm},
		arr/.style={-{Latex[length=2mm]}, thick}
		]
		\node[port] (pwr) at (0,0) {DC IN};
		\node[port, right=1.3cm of pwr] (hdmi) {HDMI};
		\node[port, right=1.3cm of hdmi] (usb3) {USB 3.0};
		\node[port, right=1.3cm of usb3] (usb2) {USB 2.0};
		\node[box, above=1.8cm of pwr] (pwrC) {5V PSU};
		\node[box, above=2.0cm of hdmi] (hdmiC) {HDMI Cable};
		\node[box, above=2.2cm of usb3] (usbC1) {USB Camera};
		\node[box, above=2.4cm of usb2] (usbC2) {Keyboard / Mouse};
		\draw[arr] (pwrC.south) -- (pwr.north);
		\draw[arr] (hdmiC.south) -- ++(0,-0.3) -| (hdmi.north);
		\draw[arr] (usbC1.south) -- ++(0,-0.3) -| (usb3.north);
		\draw[arr] (usbC2.south) -- ++(0,-0.3) -| (usb2.north);
	\end{tikzpicture}
	\caption{Cable pack to port mapping used during assembly. Each arrow indicates the correct connection for power, display, camera, and peripherals.}
\end{figure}

The cable mapping provides a simplified “port-to-cable” visualization, minimizing wiring errors during setup.  
This approach reflects the design conventions used in modern hardware documentation to improve assembly accuracy and reproducibility \cite{hardware-manuals, embedded-docs}.  
Consistent use of short, labeled cables improves electromagnetic compatibility (EMC) and reduces signal degradation, as noted by \cite{emc-wiring, cable-design}.

% -------------------------------------------------------